\chapter*{Abstract}
\addcontentsline{toc}{chapter}{Abstract}

\textbf{Tipul lucr\u arii:} sintez\u a \c si cercetare aplicativ\u a

\vspace{1em}

Lucrarea de fa\c t\u a descrie procesul de proiectare \c si construire a
aplica\c tiei web \textbf{Gym-Connect}, un instrument dedicat sportivilor
care practic\u a antrenamente de for\c t\u a \c si doresc s\u a \^ i\c si
urm\u areasc\u a progresul \^ in mod structurat. Ideea a pornit dintr-o
frustrare real\u a: majoritatea aplica\c tiilor similare de pe pia\c t\u a
sunt fie prea complicate, fie cer bani pentru func\c tii de baz\u a.
Gym-Connect \^ i\c si propune s\u a fie altfel -- gratuit, simplu \c si
accesibil direct din browser, f\u ar\u a instalare.

Din punct de vedere tehnic, proiectul urm\u are\c ste patru obiective
principale: construirea unei arhitecturi client-server func\c tionale
\^ in care frontenidul React comunic\u a cu un server Express prin cereri
HTTP autentificate; implementarea autentific\u arii utilizatorilor prin
Clerk; realizarea unui modul de \^ inregistrare a antrenamentelor \^ in timp
real, cu suport pentru calculul greut\u a\c tilor \c si indicatori de
intensitate precum RIR \c si 1RM estimat; \c si, nu \^ in ultimul r\^ and,
construirea unui modul de analiz\u a a progresului care folose\c ste
regresia liniar\u a pentru a estima evolu\c tia pe termen scurt \c si a
detecta automat platourile de performan\c t\u a.

Lucrarea este organizat\u a \^ in \c sase capitole. Primele dou\u a
stabilesc contextul -- de ce exist\u a aceast\u a problem\u a \c si ce
solu\c tii mai exist\u a pe pia\c t\u a. Urm\u atoarele dou\u a intr\u a
\^ in detaliile tehnice ale solu\c tiei propuse. Ultimele dou\u a prezint\u a
aplica\c tia concret, a\c sa cum arat\u a \c si func\c tioneaz\u a
ea, \c si trag concluziile.

\vspace{1em}
\noindent\textbf{Cuvinte cheie:} aplica\c tie web, antrenamente de
for\c t\u a, progressive overload, React, TypeScript, Express.js,
Prisma, PostgreSQL, Clerk, analiz\u a progres, 1RM, RIR.
